% !TeX root = ../papers/nc-p3.tex
\section{長さ四の基点のブローアップと十九次の排除}
\label{sec:cofactor-nineteen-clusters}
\begin{theorem}\label{thm:cofactor-exclude-nineteen}
被約窓で整な一般被約の階数二曲線の次数十九は生じない。
\end{theorem}
\begin{proof}[基点スキーム]
定理\ref{thm:cofactor-nineteen-surface}の二つのモデル $S_0=\mathbb F_0,\mathbb F_2$ を使う。
核超平面を $L$、一般母線を $F$ として $H_0=2L+F$。四つの入力切断の空間を $V$ とする。
まず、これらの共通イデアルは $S_0$ 上で長さ四の局所完全交叉 $Z$ を定めることを示す。

単純四基点の場合、順次ブローアップした滑らかな曲面 $S$ で入力束は
$H=H_0-\sum e_i$（全引戻し例外類）。各例外曲線の固有変換での次数は
一から近接点の個数を引いた数で、非負である。従って各根の上では分岐も衛星点もない自由な鎖である。
長さ $r$ の鎖の通過条件は局所座標で $(v,u^r)$ に属することである。
実際鎖を通る滑らかな枝を $v=0$ に選び、その枝に沿って $v/u,v/u^2,\dots$ のチャートで順に一点ずつ引く。
各鎖の条件を合わせて長さ四の $Z$ を取る。$V$ はこのイデアルを生成する。
局所的には最終例外直線における $H$ の次数が一で、そこでの切断は全射する。
最終チャート $v=u^r w$ で $u^r$ を除いた切断の制限は $1,w$ の線形結合で、
係数は $(u^r,v)$ の最小二生成元に関する剰余係数である。Nakayamaの補題を適用する。

二重基点の場合、通常の一点ブローアップ $S$ で既に系は自由である。
例外直線上の先頭二次式は基点なしで、二または三次元を張る。
三次元なら例外直線の入力像は滑らかな円錐曲線で、そこへの射は埋込みである。
このとき入力射と核射の対 $S\to\mathbb P^3\times K$ は閉埋込みになる。
滑らかな核曲面ではブローアップの非自明な繊維は例外直線だけで、その上で入力射が点と接方向を分離する。
法方向には核射の微分が非零である。錐の場合、二重基点は $H_0 E=1$ により負切断 $E$ の外にある。
$E$ でも入力束は次数一で、核射は法方向の微分が非零である
（頂点近くでは $w(1,t,t^2)$ の表示を使う）。従って射の対は固有、単射、非分岐である。
すると一般入力平面による正規化曲線 $B$ の各入力繊維は、核座標で閉埋込みになる。
補題\ref{lem:cofactor-plane-evaluation}の証明はこの繊維ごとの埋込みだけで足り、$g(B)=2$ を与える。
一方 $H^2=8$, $K_S H=-10+2=-8$ なので $g(B)=1$。従って先頭二次式は二次元である。
その基底を持つ二切断は互いに素な先頭二次式を持つ。局所完全交叉の長さは四、
随伴次数環は二つの先頭二次式の完全交叉で極大イデアルの三乗が零となる。
他の切断はこの二つの一次結合を引けば位数三以上なので、同じイデアルに入る。

\end{proof}
\begin{proof}[有限でない入力繊維の個数]
$\Gamma=\operatorname{Bl}_Z S_0$ とし、$H$ を $H_0$ の引戻しにイデアルの可逆引戻しを掛けた束とする。
$V$ は $H$ を生成して入力射 $h:\Gamma\to D$ を与え、$L$ は核座標を与える。
$\Gamma$ は整な局所完全交叉である。実際基点の正則列を $a,b$ とすれば
局所表示は $S_0\times\mathbb P^1$ 内の $aY-bX=0$ である。二重基点の場合には例外直線に沿う正規性を要しない。
各台点上の例外直線の $H$ 次数は一であり縮約されない。

入力に縮約される曲線は、従って $S_0$ 上の曲線の固有変換だけである。
滑らかな二次曲面で二次数 $(a,b)$ の両方が正なら各基点重複度は $\min(a,b)$ 以下
（両母線との交点数と無限近接での重複度の単調性）。
入力次数は $2b+3a-4\min(a,b)>0$ 以上である。二重基点でも減少は $2\min(a,b)$ 以下である。
次数二の母線が縮約されるには単純基点二個か二重基点が必要で、該当する母線は高々二本。
他方の次数三の母線の縮約には単純基点三個が必要で高々一本である。
$\mathbb F_2$ で母線と $E$ 以外の曲線は $aE+bF$, $a>0$, $b\ge2a$ で
各重複度は $a$ 以下。入力次数は $2b+a-4a>0$ である。母線以外は $E$ の一本だけが残る。
従って二つのモデルとも $h$ の正次元繊維は高々三つの入力点に限られる。

$M=\operatorname{coker}N\hookrightarrow h_*L$ は一般同型な評価包含である。
$h$ の有限な開集合全体ではこれは同型である。実際 $\chi(\Gamma,L+tH)=4(t+1)^2$。
これを計算するには $Z$ を相異なる四点に平滑化してよい。滑らかな曲面の点のHilbertスキームの
既約性と局所完全交叉の開性によりこの変形があり、上の一方程式表示は相対ブローアップの平坦性と
各繊維でのブローアップとの一致を示す。滑らかな四点のブローアップで
$L^2=2$, $LK=-4$, $H^2=8$, $HK=-6$, $HL=5$ にRiemann--Rochを適用する。
高次押出しは有限個の正次元繊維でのみ寄与し、$M$ も同じHilbert多項式を持つので
包含は余次元一で等号である。純次元二の $M$ は最大Cohen--Macaulayであり、有限集合を除いた切断が延長する。
両層を分数体内で比べれば有限な開集合での等号が従う。
よって二十個の階数一の印のうち少なくとも十七個は $\Gamma$ の長さ三の入力繊維である。
各々は長さ三部分スキームが一点に写る軌跡の孤立点である。有限な開集合では全繊維の長さは余階数に等しく、
階数一点は有限、長さ三の繊維中の長さ三部分スキームは一つだけだからである。
\end{proof}

\begin{proof}[基点と入力四切断の変形]
中央の $V$ を四相異点の一般の四切断へ変形できる。単純鎖の場合、滑らかな解消 $S$ で
$K_SH=-6$、一般入力切断 $B$ は種数二で $H|_B$ の次数八である。
$S$ は有理なので切断完全列から $h^0(S,H)=8$。これは上記の鎖の記述で
$h^0(S_0,I_Z H_0)$ に等しい。二重基点では $H_0$ はそこで二次ジェットを全て与える。
積二次曲面では単項式から明らかであり、$\mathbb F_2$ の負切断外の局所座標なら
$A_5(t)+B_3(t)z+C_1(t)z^2$（添字は次数上限）が全二次ジェットを与える。
どちらも $h^0(H_0)=12$, $H^1(H_0)=0$（母線方向の押出しで計算）なので
$H^1(I_Z H_0)=0$。従って $Z$ と曲面を動かす近傍で八次元切断束を得る。
$\mathbb F_2$ から $\mathbb F_0$ へも移れる。$\mathcal O(2)$ の $\mathcal O$ による拡大を
分裂点から非分裂点へ動かし、射影束の相対超平面束 $L$ と母線束 $F$ を使う。
滑らかな曲面族の四点の相対Hilbertスキームは滑らかで各繊維が既約であり、
局所完全交叉の開集合上で上述の平坦なブローアップを作れる。八次元切断束の四空間とその基底を動かせばよい。

一般繊維の長さ三縮約スキームが有限となることも示す。$\mathbb F_0$ の相異なる四母線に一つずつ中心を置く。
母線への直像 $W$ は一般繊維で階数三であり、元の $\mathcal O(3)^3$ に四つの値条件を課した束になる。
各特別繊維は交わる二直線で入力次数は $(1,1)$、一般繊維は次数二の滑らかな有理曲線である。
従って相対完全系で滑らかなブローアップは $\mathbb P(W)$ 内に繊維ごとに平面二次曲線として埋まる。
四中心を他方の母線座標 $0,0,\infty,\infty$ に取れば
$W=\mathcal O(1)\oplus\mathcal O(3)\oplus\mathcal O(1)$ であり、一般にも全直和次数が正となる。
完全系はこの三次元有理正規スクロールを $\mathbb P^7$ に埋める。
曲面の類は $2\xi-2F$ である（相対次数二、$\xi^3=5$ と $\xi^2[S]=8$）。
スクロールは二次式で定義され、射影正規である。曲面方程式に底の二次式を全て掛けると
スクロール上の二次超平面切断になり、これらは射影空間の二次式に持ち上がる。
従って長さ三で曲面に会う直線は曲面内に含まれる。

一般の四切断の射影中心は曲面を避けるので、曲面内の直線を潰さない。
完全系で平面を張る長さ三スキームに対しては、四切断の評価が階数一以下となる条件は余次元六である。
滑らかな曲面の三点のHilbertスキームの次元は六であり、一般の四空間の退化軌跡は有限である。
長さ三が直線を張る場合は前段落で既に除かれた。
\end{proof}

\begin{proof}[孤立点の持ち上げと次数十六]
相対ブローアップの三点のHilbertスキームで、四切断から階数三の同伴束 $H^{[3]}$ への評価の
階数一以下を取る。中心繊維の各孤立点は一般繊維に移る。実際各長さ三スキームは
埋込み次元二以下で、滑らかな $r$ 次元の局所埋込空間の三点のHilbertスキームはそこで滑らかな次元 $3r$
（一般平面射影に関するグラフの表示を使う）。相対曲面は局所完全交叉なので
近傍で長さ三をその上に置く条件は高々 $3(r-2)$ 個の方程式である。分離した台では長さを分配する。
四切断は基点なしで評価の階数は少なくとも一なので、階数一となる条件はさらに六式で書ける。
一般点へ向かう一次元基底を含めて局所次元は一以上であり、中心で孤立なら水平成分が必要である。
相異なる中心孤立点は一般繊維の長さに個別に寄与する。例えば固有な相対退化軌跡を
hensel化した底上で考え、特殊繊維の各孤立点を連結成分として分ければよい。

一般繊維の長さはThom--Porteous公式で
\[
 \int_{S^{[3]}}(c_2^3-2c_1c_2c_3+c_3^2)=16,\qquad c_i=c_i(H^{[3]}).
\]
積分の再現手順を記す。滑らかな四点ブローアップの族で積分値は一定なので四中心を積二次曲面の
トーラス固定四隅に置く。各固定点で接重みを $u,v$、線束重みを $h$ とすると、
一点ブローアップ後の二組は $(u,v-u,h-u)$, $(u-v,v,h-v)$。
元の四隅 $(i,j)\in\{0,1\}^2$ では $((1-2i)U,(1-2j)V,-2iU-3jV)$ を使う。
各固定点に合計サイズ三の分割を割り当て、行 $i$ 列 $j$（零始まり）の箱は同伴束に
$h-iu-jv$ を与える。分母は各箱の腕長 $a$、脚長 $l$ について
$((a+1)v-lu)((l+1)u-av)$ の積である。これは滑らかな曲面の点のHilbertスキームの
単項式固定点での接重み公式である。三つの同伴束重みの基本対称式を上の分子に代入し、
全固定点で割り算して加える。\path{computations/verify_cofactor_nineteen_clusters.py}
は二組の整数重みでこの有理数和を検算する。十七個以上の孤立点に対して長さ十六は不可能である。
\end{proof}
